\section{Métodos}

Esta seção detalha a arquitetura, a infraestrutura e a metodologia adotada para o desenvolvimento do sistema de recomendação acadêmica. O fluxo do projeto é estruturado em três frentes principais: a interface de interação com o discente via chatbot, a extração e o processamento dos dados históricos, e a aplicação do modelo de agrupamento para a geração de recomendações.

\subsection{Análise da base de dados}

O projeto previa a utilização de históricos escolares reais de discentes do curso de Ciência da Computação da Universidade Federal Rural do Rio de Janeiro (UFRRJ) como fonte primária de dados. Entretanto, não foi possível utilizar a base de dados real, devido à dificuldade de acesso aos dados dos alunos, embora esses dados tenham sido formalmente solicitados à Coordenadoria de Tecnologia da Informação e Comunicação (COTIC) da universidade.

Para verificar a viabilidade da execução do projeto, foram criados dados hipotéticos, que permitiram testar o pipeline de pré-processamento e a aplicação dos algoritmos de agrupamento e de classificação. Esses dados foram gerados a partir de uma pequena amostra real de hist\'oricos acad\^emicos, reproduzindo estatisticamente as distribui\c{c}\~oes de notas e situa\c{c}\~oes de aprova\c{c}\~ao/reprova\c{c}\~ao observadas nos dados originais. Ressalta-se que esses experimentos t\^em car\'ater metodol\'ogico e que seus resultados n\~ao representam a comunidade discente do curso.

\subsection{Análise da viabilidade de LLMs}

A fim de prover uma interface conversacional e personalizada, avaliou-se a incorporação de Grandes Modelos de Linguagem (Large Language Models, LLMs) à arquitetura do tutor inteligente. A integração desses modelos visava permitir uma interação mais natural com os estudantes, apresentando as inferências geradas pelos algoritmos de Mineração de Dados Educacionais de forma contextualizada e adaptada às necessidades de cada discente.

Para tanto, analisaram-se diferentes modelos e serviços de LLM disponibilizados por provedores como OpenAI, Google e Meta. A investigação focou na interoperabilidade dessas ferramentas com a arquitetura proposta, observando sua integração com os dados acadêmicos estruturados e com os algoritmos de aprendizado de máquina empregados no sistema.

A análise de viabilidade técnica e econômica considerou o custo de consumo das APIs, os limites de tokens, a vazão de requisições por unidade de tempo, o suporte a requisições simultâneas e a escalabilidade no atendimento ao público-alvo. Entretanto, essas ferramentas são pagas, o que não seria o ideal para um projeto de iniciação científica desenvolvido sem aporte financeiro, inviabilizando o uso de APIs de modelos proprietários.

Embora alguns provedores ofereçam níveis gratuitos (\textit{free tiers}), os tetos de requisições observados mostraram-se insuficientes para sustentar a demanda de múltiplos acessos simultâneos. Essa restrição comprometeria a disponibilidade e o tempo de resposta do sistema, inviabilizando sua aplicação no curso de Ciência da Computação e sua eventual expansão para demais cursos da universidade.

Diante dessas restrições orçamentárias e operacionais, a incorporação de LLMs à arquitetura atual foi postergada. Optou-se por priorizar o desenvolvimento do tutor inteligente fundamentado estritamente nos algoritmos analíticos e no fornecimento de informações acadêmicas objetivas (como período, matérias, orientações sobre o fluxo do curso e datas importantes do calendário acadêmico), mantendo a integração com modelos de linguagem como uma perspectiva de trabalhos futuros, condicionada à alocação de recursos financeiros ou computacionais adequados.

\subsection{Interação com o Usuário e Coleta de Dados via ChatBot}
Como ponto de contato principal com os estudantes, o sistema foi concebido como um assistente virtual (ChatBot) integrado à plataforma Telegram via API, escolha motivada pela ampla utilização da ferramenta por parte dos estudantes, pela alta acessibilidade que proporciona e pelo acesso gratuito à sua API, que permite a integração sem custos adicionais. Além de servir como canal para o envio do documento de histórico acadêmico, arquivo fundamental para a extração dos dados brutos de notas, o bot foi desenvolvido para atuar como uma ferramenta completa de suporte acadêmico.

Além da coleta de dados, o sistema disponibiliza funcionalidades voltadas à consulta e à orientação acadêmica dos discentes. Nesse sentido, buscou-se oferecer um suporte abrangente, reunindo em uma única interface informações que, de outro modo, estariam dispersas em diferentes canais institucionais.

Dentre suas funcionalidades, a aplicação oferece:
\begin{itemize} 
    \item \textbf{Dicas e Orientações:} sugestões estratégicas sobre o fluxo do curso;
    \item \textbf{Consulta ao Calendário Acadêmico:} consulta às informações do calendário vigente no período letivo, como as datas importantes (prazos de cadastro de turmas, pré-matrícula, ajuste de grade, rematrícula e encerramento do semestre).
\end{itemize}

\subsection{Processamento do Histórico Acadêmico e Agrupamento com K-Means}
A partir do momento em que o aluno submete o seu histórico acadêmico através do ChatBot, o sistema realiza a leitura e a extração estruturada das notas obtidas em cada disciplina cursada. Esses dados brutos alimentam a etapa de aprendizado de máquina, que emprega o algoritmo não supervisionado K-Means \citep{huberty2005cluster}, implementado por meio da biblioteca scikit-learn \citep{pedregosa2011scikit}.

O K-Means opera analisando o vetor de notas do discente e comparando-o com o padrão de desempenho de outros estudantes. O objetivo desta segmentação por agrupamentos (clusters) é identificar semelhanças de trajeto e reconhecer quais disciplinas funcionam como gargalos ou apresentam maiores índices de dificuldade em comum para determinados perfis de alunos, automatizando o mapeamento de vulnerabilidades no fluxo do curso.

\subsection{Classificação com KNN}

Como etapa complementar à análise de perfis realizada pelo K-Means, aplicou-se o algoritmo K-Nearest Neighbors (KNN) para realizar duas formas de classificação dos estudantes.

A primeira classificação teve como objetivo identificar o eixo de formação (Matemática, Programação ou Sistemas) no qual cada estudante apresentou maior desempenho acadêmico. A escolha desses três eixos, associados às disciplinas dos primeiros períodos do curso, foi motivada pela maior disponibilidade de dados referentes às disciplinas introdutórias e pela relevância dessas etapas iniciais para a progressão acadêmica dos estudantes. Embora o curso contemple outros eixos de formação ao longo dos períodos subsequentes, estes não foram incorporados à análise nesta etapa devido à menor disponibilidade de amostras.

Como características de entrada do classificador por eixo, foram utilizadas as notas individuais das disciplinas pertencentes aos três eixos, totalizando onze variáveis: quatro disciplinas de Sistemas, três de Matemática e quatro de Programação. Antes da aplicação do KNN, os dados foram padronizados por meio do método StandardScaler, de modo a colocar as diferentes variáveis em uma escala comparável. Para cada estudante, calcularam-se as médias das disciplinas pertencentes a cada grupo disciplinar, atribuindo-se a classe correspondente ao eixo de maior média. Essa classificação foi realizada com k = 5, considerando os cinco vizinhos mais próximos para determinar a classe de cada estudante.

Em uma segunda aplicação, o KNN foi utilizado para classificar os estudantes de acordo com o desempenho acadêmico geral, distribuindo-os em quatro faixas definidas a partir da m\'edia geral:
[0, 2{,}5); [2{,}5, 5{,}0); [5{,}0, 7{,}5); e [7{,}5, 10{,}0]. Essa classificação permite representar de maneira mais abrangente o desempenho acadêmico do estudante, complementando a identificação do eixo de maior desempenho.

Em conjunto, essas duas classificações possibilitam caracterizar tanto a área de maior afinidade acadêmica de cada estudante quanto sua situação de desempenho geral, informações que, associadas a dados sobre disciplinas pendentes, pré-requisitos e histórico, poderão futuramente apoiar a construção de recomendações personalizadas de trajetória acadêmica.

\subsection{Sistema de Recomendação de Fluxo Acadêmico}
No estágio atual de desenvolvimento da ferramenta, o foco metodológico concentra-se na análise profunda do componente curricular e no mapeamento e identificação do perfil discente. A partir dos agrupamentos gerados pelo K-Means, o sistema é capaz de posicionar o estudante em um \emph{cluster} que reflete suas semelhanças de desempenho em relação a pares e históricos anteriores. 

Embora a recomendação automatizada de disciplinas específicas utilizando Inteligência Artificial constitua a etapa final da arquitetura proposta -- em linha com
abordagens como a de \citet{wagner2024courserecommender}, que recomendam disciplinas
com base nas trajet\'orias de estudantes bem-sucedidos -- esta
funcionalidade encontra-se em fase de estrutura\c{c}\~ao.